\title{LongRoPE: Extending LLM Context Window Beyond 2 Million Tokens}

\begin{abstract}
Having a large context window is an important capability for large language models (LLMs). Nevertheless, existing extended context windows typically max out at approximately 128k tokens, owing to substantial costs associated with fine-tuning, the limited availability of sufficiently lengthy textual data, and the detrimental impact of introducing new token positions that often lead to instability.
\end{abstract}

This work presents LongRoPE, a novel approach that, for the first time, pushes the context window of pre-trained LLMs up to an unprecedented \textbf{2048k} tokens, requiring at most 1k fine-tuning iterations conducted with context lengths of up to 256k tokens—all while preserving model performance within the standard short context window. The method builds upon three principal innovations: (i) by systematically identifying and leveraging two types of non-uniformities in positional interpolation through an efficient search process, we achieve superior fine-tuning initialization and permit up to 8$\times$ context expansion without the need for fine-tuning; (ii) a staged extension procedure is proposed, whereby we first adapt an LLM to 256k token sequences, then apply a secondary round of positional interpolation to this fine-tuned model to further enlarge the context window to 2048k; (iii) we recalibrate LongRoPE at the 8k-token context length to fully restore original short-window performance. Comprehensive experiments carried out on both LLaMA2 and Mistral over a suite of benchmarks substantiate the benefits of our methodology. The models enhanced via LongRoPE preserve their foundational architecture, altering only the positional embedding component, and thus remain compatible with most existing optimization strategies. Source code will be provided at \texttt{https://github.com/microsoft/LongRoPE}

\section{Introduction}

Although Large Language Models (LLMs) have achieved significant advancements across a range of applications~\cite{gpt4,llama2}, they typically face restrictions imposed by a finite context window, such as the 4096-token constraint of LLaMA2~\cite{llama2}. When the input exceeds this boundary, the models' effectiveness diminishes, primarily because they have not been exposed to such extended positional contexts during training. This limitation creates obstacles in key use cases, including in-context learning that demands a large set of exemplars~\cite{Huang2023FewerIM}, as well as the deployment of LLM agents~\cite{park2023generative,madaan2023selfrefine}.

Recent studies indicate that the context window of a pre-trained LLM can be broadened to approximately 128k through fine-tuning with longer input sequences~\cite{longlora,pi,yarn,zhang2024soaring,liu2023scaling}. Nevertheless, scaling the context window further encounters three primary challenges. Firstly, the introduction of previously unseen positional indices during extension often results in catastrophic values, which in turn cause distributional shifts and hinder successful fine-tuning convergence~\cite{pi}. This issue becomes significantly pronounced when the context is expanded from 4k to over 1000k, as more than 90\% of the positions represent new, untrained indices. Secondly, effective fine-tuning generally relies on the availability of sequences that match the target length. However, current datasets contain very few sequences of such extreme length, particularly those beyond 1000k, and training on these extended texts requires substantial computational resources, including extensive time and GPU availability. Thirdly, scaling up to extremely large context windows causes the model’s attention mechanism to become overly diluted across many token positions, ultimately reducing effectiveness for tasks involving the original, shorter context sizes~\cite{pi}.

A commonly used strategy to address the initial challenge involves interpolating the RoPE positional embedding~\cite{rope,pi}, where new positional indices are compressed into the range encountered during pretraining, as illustrated in Fig.\ref{fig:overview}. Position Interpolation (PI)~\cite{pi} achieves this by linearly adjusting RoPE's rotary angles according to the scaling factor. In contrast, NTK~\cite{ntk,dynamicntk} employs non-uniform interpolation and extrapolation techniques that differ across the RoPE dimension space. YaRN~\cite{yarn} takes a frequency-oriented approach, segmenting RoPE dimensions into three separate frequency bands and assigning extrapolation, NTK-based, or linear interpolation methods to each group. Despite these innovations, Transformer models' positional embeddings inherently possess \emph{intricate, non-uniform information entropy}. Present methods do not sufficiently exploit this complex variability, which results in loss of information and constrains the maximum usable context window.

Section~\ref{sec:analysis} provides empirical evidence for two main observations: \textbf{(1)} Successful positional interpolation necessitates accounting for two distinct non-uniform factors: variability across RoPE dimensions and across token positions. It is observed that lower-dimensional RoPE components and tokens appearing earlier in the sequence generally require less interpolation; however, the optimal interpolation strategy must be tailored to the desired sequence extension length.
  \textbf{(2)} Incorporating these identified non-uniformities within the positional interpolation approach enables more effective preservation of the original RoPE’s information, particularly in the most critical dimensions and early token locations. This approach curtails the information loss due to interpolation, resulting in superior initial parameters for subsequent fine-tuning. Furthermore, this method enables up to an 8$\times$ length extension even when fine-tuning is not conducted.

Building upon these observations, we propose \textbf{LongRoPE}, a robust technique that pushes the boundaries of the LLM context window to surpass 2 \emph{million} tokens. The design of LongRoPE relies on three principal advancements. The first involves harnessing the multidimensional non-uniformity present in positional interpolation to its fullest potential. Specifically, {LongRoPE} determines optimized rescale factors for the rotation angles in each dimension of RoPE by taking into account specific token positions. Given that the search space for these rescale factors grows exponentially as the desired extension ratio increases, {LongRoPE} incorporates an evolutionary search algorithm complemented by two optimization strategies to enhance the efficiency of this process. An illustration of the optimized rescaled RoPE discovered through this approach is provided in Fig.~\ref{fig:overview}.

Subsequently, {LongRoPE} employs a resource-efficient, incremental expansion approach to attain a 2048k context window, circumventing the need for fine-tuning on ultra-long sequences, which are both rare and difficult to obtain. The approach initiates by identifying a 256k context length within the pre-trained LLM and performing fine-tuning at this range. Thanks to our non-uniform positional interpolation, which enables up to an 8$\times$ extension even without fine-tuning, we carry out a follow-up search to determine new RoPE rescaling parameters for the already fine-tuned, length-augmented LLM. Through this method, a 2048k context window is realized for both LLaMA2 and Mistral~\cite{mistral}.

In order to address potential performance declines when using the original, shorter context windows, {LongRoPE} applies further modifications to the RoPE rescale factors in the expanded LLM. Analogous to the upscaling process from 256k to 2048k, the model is adjusted down to 4k and 8k context lengths on the LLM fine-tuned at 256k by leveraging our search algorithm, aiming to minimize positional interpolation. For inference scenarios where the sequence length falls below 8k, RoPE is modified to employ the rescale factors identified by this search.

Comprehensive evaluations conducted on multiple LLMs and a range of long-context tasks validate the efficacy of our approach. Our results indicate that LongRoPE consistently sustains low perplexity across evaluation lengths spanning from 4k to 2048k, attains more than 90\% passkey retrieval accuracy, and matches the accuracy of established benchmarks created within a 4096 context window. LongRoPE is compatible with all LLMs that utilize RoPE embedding. We plan to make both our code and the LongRoPE-2048k models publicly available.

\section{Non-uniformity in Positional Interpolation}
\label{sec:analysis}

\subsection{Preliminary}

Transformer models require explicit positional information, often in the form of position embedding, to represent the order of input tokens. Our work focuses on the RoPE~\cite{rope} position embedding, which is widely used in recent  LLMs. For a token at position index $n$, its corresponding RoPE encoding can be simplified as follows:

\begin{equation}
		\fontsize{7.8}{7.8} \selectfont
	\label{eq:rope}
	\left[ cos(n\theta_0), sin(n\theta_0), cos(n\theta_1), \cdot\cdot\cdot, cos(n\theta_{d/2-1}), sin(n\theta_{d/2-1}) \right]   
\end{equation}

Here, $d$ indicates the dimensionality of the embeddings, $n\theta_i$ denotes the rotary angle assigned to the token situated at position $n$, and $\theta_i = \theta^{-2i/d}$ specifies the frequencies used for rotation. In the RoPE mechanism, the typical value chosen for $\theta$ is 10000.

\textbf{\emph{Context window extension ratio $s$} and positional interpolation}. The parameter $s$ is defined as the quotient of the new, extended context window length $L'$ by the original window length $L$, formally expressed as $s = \frac{L'}{L}$.

To extend the context window from $L$ to $L'$, current positional interpolation methods suggest downscaling rotation frequencies $\theta_i$ based on extension ratio $s$. Let $\beta=\theta^{2/d}$, and $\lambda$ denote the actual rescale factor related to $s$, we   unify these positional interpolation methods as follows:

\begin{equation}	
	\fontsize{7.5}{7.5} \selectfont
	\label{eq:unified_rope}
	\begin{aligned}
		\left[cos\left(\frac{n}{\lambda(\beta)^0}\right), sin \left(\frac{n}{\lambda(\beta)^0}\right), cos\left(\frac{n}{\lambda(\beta)^1}\right), \cdot\cdot\cdot,  sin \left(\frac{n}{\lambda(\beta)^{d/2-1}}\right) \right]   
	\end{aligned}
\end{equation}

\textbf{Linear positional interpolation (PI)}. PI~\cite{pi} introduces a technique where position indices are linearly interpolated when remaining within the maximum sequence length seen during pre-training. Specifically, for an intended extension ratio $s$, the rotational angles assigned to each position are uniformly scaled down by a factor of $\lambda = s$ in every RoPE dimension. This uniform compression causes positional representations to become densely packed, which in turn impairs the model's capacity to separate tokens that are near each other in the sequence. As a consequence, the effectiveness of PI diminishes notably when the extension ratio is large.

\textbf{NTK-based interpolation and extrapolation}. ~\cite{ntk,dynamicntk} reinterpret RoPE as a form of information encoding and utilize the Neural Tangent Kernel (NTK) framework~\cite{jacot2018neural,tancik2020fourier}. To address the problem of positional crowding inherent in PI, they propose redistributing the interpolation constraint among different RoPE dimensions. Specifically, they decrease the scaling for lower (corresponding to high-frequency) dimensions and increase it for higher (associated with low-frequency) dimensions. This adjustment enables the model to perform both positional interpolation and extrapolation using $\lambda=s^{i}$. The dynamic NTK variant~\cite{dynamicntk} further modifies the extension ratio at every position as a function of the current sequence length. Contrary to PI approaches, which require fine-tuning, NTK-informed strategies can expand context windows even without additional fine-tuning; however, practical extension is typically capped at around 4$\times$.

\textbf{YaRN}~\cite{yarn} presents a notable advancement in positional interpolation by organizing RoPE dimensions into three distinct groups based on their frequency components, assigning a separate interpolation method to each group. Specifically, high frequency dimensions are processed with extrapolation ($\lambda$=1), low frequency dimensions are interpolated linearly (PI), and those in the intermediate range use NTK. The central aspect of YaRN is this frequency-based partitioning of RoPE dimensions, which, however, is currently derived from manual empirical tuning. Consequently, this reliance on human heuristics may limit its effectiveness when applied to novel LLM architectures.

\subsection{Study on Non-uniform Positional Interpolation}

Drawing motivation from NTK and YaRN, we observe that their improvements stem from incorporating non-linearities, particularly through the selective treatment of various frequencies along RoPE dimensions to facilitate tailored interpolation and extrapolation. Nevertheless, present approaches to non-linearity depend extensively on manually crafted heuristics. This situation leads us to two central inquiries: (1) Can the effectiveness of existing positional interpolation methods be improved? (2) Might there exist additional non-linear strategies that remain untapped?

To address these questions, we employ evolution search (refer to Sec.~\ref{sec:method}) to identify improved non-uniform positional interpolation strategies for LLaMA2-7B. This search process is directed by perplexity scores, utilizing 5 randomly selected samples from the PG19~\cite{pg19} validation set. Our empirical study yields the following important observations.

\textbf{Finding 1}: \textit{RoPE dimensions exhibit substantial non-uniformities, which are not effectively handled by current positional interpolation methods.}

We optimize $\lambda$ for each RoPE dimension as specified in Eq.~\ref{eq:unified_rope}. In Table~\ref{tbl:exp1}, we report the perplexity results of LLaMA2-7B evaluated on PG19 and Proof-pile~\cite{proof-pile} test data using various methods, all in a zero-shot setting without model fine-tuning. Our optimized approach consistently yields notable reductions in perplexity, indicating that widely-used interpolation strategies—both linear (PI) and more complex forms (Dynamic-NTK and YaRN)—do not reach the best performance possible. Interestingly, on PG19, YaRN results in worse perplexity than both PI and NTK, attributable to its inability to extend to the maximum sequence length for LLMs that have not been fine-tuned. As an illustration, when tested with an 8k context window, YaRN’s perplexity rises substantially once the sequence surpasses 7k tokens.

In our investigation, the rescaled factors $\lambda$ described in Eq.~\ref{eq:unified_rope} are not uniform across dimensions, unlike the constant scale $s$ used in PI, NTK’s approach, or the group-based computation in YaRN. This introduction of non-uniform rescaling leads to notable improvements in LLaMA2’s language modeling, as reflected in reduced perplexity for context windows of 8k and 16k, even in the absence of any additional fine-tuning. The improvement is attributed to the modified positional embedding, which largely retains the integrity of the original RoPE—particularly along essential dimensions—thereby easing the model’s challenge of discriminating between proximate token positions.

\textbf{Finding 2}: \textit{RoPE for the initial tokens in the input sequence should be extrapolated with less interpolation.} 

We propose that position interpolation for the initial $\hat{n}$ tokens in input sequences should be reduced, since these early tokens tend to receive higher attention scores and are particularly influential within attention layers, as previously reported in Streaming LLM~\cite{streamingllm} and LM-Infinite~\cite{han2023lminfinite}. To assess this, we conduct experiments by expanding the context window to 8k and 16k using both PI and NTK approaches, selectively omitting interpolation for the first $\hat n$ (0, 2, ..., 256) tokens. The case $\hat n=0$ corresponds to the original PI and NTK. Results summarized in Table~\ref{tbl:tokenposition} reveal two principal findings: \textbf{(1)} bypassing position interpolation for the initial tokens leads to noticeable performance gains, and \textbf{(2)} the value of $\hat n$ that yields optimal results is contingent upon the intended length of the extended context window.

\textbf{Finding 3}: \textit{Non-uniform positional interpolation effectively extends LLM context window in both fine-tuning and non-fine-tuning settings.} 

Our findings demonstrate that the non-uniform position interpolation we devised considerably boosts extension performance for 8k and 16k context lengths, even in the absence of fine-tuning. However, when scaling to even larger context windows, such as 64k, fine-tuning becomes necessary. Therefore, we apply fine-tuning to LLaMA2-7B using our tailored RoPE adaptation for a 64k window (refer to the Appendix for details on the fine-tuning procedure). As evidenced by the results in Table~\ref{tbl:finetune}, our approach achieves substantial gains over both PI and YaRN, whether fine-tuning is applied or not. This improvement can be attributed to our strategic employment of non-uniform positional interpolation, which reduces information degradation and yields superior initial conditions for subsequent fine-tuning.

\textbf{Summary}. Our research highlights two distinct forms of non-uniformity—across RoPE dimensions and token positions. By leveraging these aspects in our positional interpolation strategy, we can markedly enhance LLM context extension capabilities.

\section{LongRoPE}

\label{sec:method}
Building upon these observations, we propose LongRoPE, which initially incorporates a streamlined search algorithm designed to leverage the two distinct non-uniformities. This approach is subsequently utilized to expand the LLM context window to exceed 2 million tokens.

\subsection{Problem Formulation}

These two types of non-uniformities greatly expand the solution space, thereby increasing the difficulty of the optimization process. To tackle this challenge, we recast the optimization of multidimensional non-uniform position interpolation as a search problem.

For a LLM targeting a  context window size of $L'$ and lengthy input documents $\mathbf{X}$, where each $\mathbf{x}\in\mathbf{X}$ surpasses $L'$ in token length, we denote the original rotary angle of the $i^{th}$ dimension in RoPE embedding at token position $n$ as  $\frac{n}{\beta^i}$. The optimization problem is then formulated as follows:

\begin{equation}
\begin{gathered}
		\label{eq:problem}

\underset{\mathbf{x} \in \mathbf{X}; \, |\mathbf{x}| \geq L'}{\operatorname{arg\,min}} \, \mathcal{L} \left( \text{LLM}(\text{RoPE}, \mathbf{X}) \right), \; \text{where } 
\\
\scalebox{0.93}{
$\underset{
\begin{subarray}{l}
i=0,\ldots,\frac{d}{2}-1; \\
n \in [0, |\mathbf{x}|); 
\end{subarray}
}{\text{RoPE}_ {(n)}}
= \left[ \cdots, \cos\left(\mathbb{I}(\hat{\lambda}_{i}, \hat{n}) \cdot \frac{n}{\beta^{i}}\right), \sin\left(\mathbb{I}(\hat{\lambda}_{i}, \hat{n}) \cdot \frac{n}{\beta^{i}}\right), \cdots \right]$}\\
\text{with } \mathbb{I}(\hat{\lambda}_{i}, \hat{n}) = \begin{cases}
1 & \text{if } n < \hat{n} \\
{\frac{1}{\lambda}_{i}} & \text{if } n \geq \hat{n}
\end{cases}
\end{gathered}
\end{equation}

In this context, we define a group of scaling factors, $\mathbb{I}(\hat{\lambda}_{i},\hat{n})$, designed to accommodate both types of non-uniformities. Here, $\hat{\lambda_i}$ reflects the non-uniformity associated with the RoPE dimensions, while $\hat{n}$ corresponds to variations across token positions. The scaling factor $\mathbb{I}(\hat{\lambda}_{i},\hat{n})$ is utilized to adjust the rotational angle for the $i^{th}$ RoPE dimension, applying the scaling parameter $\hat\lambda_{i}$ beyond a certain token position threshold $\hat{n}$. For the initial $\hat{n}$-1 tokens, the rotation angle remains unmodified, taking the original form $\frac{n}{\beta^i}$. When the token position satisfies $n\ge\hat{n}$, the appropriate scaling factor is introduced.

With a desired context window length of $L'$, our aim is to determine the best rescaling coefficients ($\mathbb{I}(\hat{\lambda}_{0},\hat{n})$, $\mathbb{I}(\hat{\lambda}_{1},\hat{n})$, ... $\mathbb{I}(\hat{\lambda}_{i},\hat{n})$, and so forth) for each RoPE dimension from the $1^{st}$ up to the $d^{th}$. The intention is that, after applying these rescaled $\text{RoPE}$ parameters to the target $\text{LLM}$, the model will attain the lowest possible loss $\mathcal{L}$ for next-token predictions—that is, achieve minimal perplexity—across input sequences $\mathbf{X}$ that are $L'$ tokens in length.

\subsection{Searching the Non-uniform Position Interpolation}

In addressing the challenge outlined in Eq.~\ref{eq:problem}, we present a straightforward but powerful approach that identifies the best RoPE rescale parameters. This method is designed to make full use of the complex, multidimensional disparities inherent in position embedding.

\textbf{Search space}. Our approach constructs an expansive search space to accommodate both forms of non-uniformity. The detailed configuration of this space is depicted in Table~\ref{tbl:searchspace}. Rather than searching over $\mathbb{I}(\hat{\lambda}_{i},\hat{n})$, we instead search for the parameters $\lambda_i$ and $\hat{n}$ directly, forgoing the inverse transformation $\hat{\lambda}_{i} = 1/\lambda_i$ to streamline the design. For each dimension in the RoPE embedding, a unique rescale factor is optimized. As specified in Table~\ref{tbl:searchspace}, $\lambda_i$ can take values ranging from 1.0 (corresponding to direct extrapolation) up to $s\times1.25$ (enabling broader interpolation than standard PI), with increments of 0.01; here, $s$ stands for the desired context window scaling factor.

The parameter $\hat{n}$ specifies how many of the starting token positions utilize the original RoPE embedding rather than applying position interpolation. In practice, we let $\hat{n}$ vary over the set \{0, 1, 2, 4, 8, 12, 16, 20, 24, 28, 32, 64, 128, 256\} during the search. Setting $\hat{n}=0$ implies that every token position employs the rescale factors determined in the search process.

\textbf{Evolution-based search}. The search space outlined in Table~\ref{tbl:searchspace} encompasses a wide array of positional interpolation options, making thorough exploration computationally demanding. As an illustration, an extension of $s=4\times$ results in $400^{128/2}\times14$=4$\times10^{167}$ different configurations. This space grows exponentially with higher extension ratios. To efficiently navigate this vast landscape, we adopt an evolution search strategy~\cite{guo2020single} and integrate two optimization strategies to markedly enhance the speed of the search. The complete search process is detailed in Algorithm~\ref{alg:search}.

\textit{Optimized initial population generation}. In place of purely random initialization for the population of $P$ rescale factors, we explicitly include the three RoPE rescale factors associated with PI, NTK, and YaRN as distinct individuals within the starting population. For the other $P$-3 members, these three rescale factors are subjected to random mutation, each with a probability of $p$.

\textit{Monotonically non-decreasing constraint}. Once the initial population has been created, we evaluate the LLM perplexity for each member. This involves adjusting the RoPE rescale factors within the target LLM and determining the perplexity of the input $\mathbf{X}$. The k highest-ranking individuals are then selected as parents for the next stage of evolution. Nonetheless, due to the immense size of the search space, straightforward applications of mutation and crossover often yield suboptimal candidates, which in turn trigger numerous unnecessary perplexity evaluations. This inefficiency is exacerbated for large $L'$, as each perplexity calculation is computationally intensive and time-consuming.

To tackle this issue, a monotonicity constraint is enforced such that the sequence of sampled RoPE rescaling factors is non-decreasing: $\lambda_i\le \lambda_{i+1}$. We restrict our RoPE application in LLM perplexity evaluations to only those configurations that fulfill this constraint, which substantially narrows down the search space. More precisely, we stipulate that $\lambda_i$ must monotonically increase as a function of the RoPE dimension (where $i = 0, \ldots, 63$). This requirement for dimension-wise monotonicity is motivated by insights from NTK theory~\cite{jacot2018neural,tancik2020fourier,ntk}, which indicates that dimensions associated with higher frequencies (i.e., lower indices) demand less interpolation—thus a lower $\lambda_i$—whereas dimensions with lower frequencies (higher indices) can tolerate greater interpolation, corresponding to a higher $\lambda_i$.

\begin{algorithm}[t]
	\small
	\caption{The search algorithm}
	\textbf{Input:} target LLM, input samples $\mathbf{X}$,   population size $P$, mutation size $N_1$, crossover size $N_2$, max iterations $\mathcal{T}$, mutate probability $p$\\

	\begin{algorithmic}[1]
		\label{alg:search}
		\STATE Initialize Top-k as an empty set: Top-k=$\phi$;	
		\STATE Generate initial population using optimization: $\text{P}_0$ = \textit{Initialize\_population\_with\_optimization} ($P$, $\mathbf{X}$, $p$);
		\FOR{$i$ from $1$ to $\mathcal{T}$}
		\STATE Evaluate the perplexity for $\text{P}_{i-1}$ on the LLM using $\mathbf{X}$: \textit{Compute\_perplexity} (LLM, $\text{P}_{i-1}$, $\mathbf{X}$);
		\STATE Refresh Top-k by incorporating new candidates: Top-k = \textit{Update\_Topk} (Top-k, $\text{P}_{i-1}$);
		\STATE Derive $\text{P}_{mutation}$ via constrained mutation on Top-k: $\text{P}_{mutation}$ = \textit{Mutation\_with\_mono\_constraint} (Top-k, $N_1$, $p$);
		\STATE Construct $\text{P}_{crossover}$ through constrained crossover with Top-k: $\text{P}_{crossover}$ = \textit{Crossover\_with\_mono\_constraint} (Top-k, $N_2$);
		\STATE Form next generation: $\text{P}_i$ = $\text{P}_{mutation}$ $\cup$ $\text{P}_{crossover}$ $\cup$ Top-k;
		\ENDFOR
		\STATE Output the Top-k member exhibiting the minimum perplexity;
	\end{algorithmic}
\end{algorithm}

\textbf{8$\times$ extension without fine-tuning}. By utilizing evolutionary search, our approach discovers optimal non-uniform RoPE rescale factors that maintain important dimensions and positions, thus reducing the risk of information loss during interpolation. As illustrated in Fig.\ref{fig:non-ft}, this strategy enables the context window of LLaMA2 to be enlarged from 4k to 32k tokens without necessitating fine-tuning. By comparison, alternative techniques—such as PI, along with non-uniform NTK and YaRN—result in significantly higher perplexity beyond a 2$\times$ extension.

\subsection{Extending LLM Context Window to 2048K}
\label{sec:extendto2048k}

\textbf{Progressive extension to 2048k}. In this section, we present our approach for expanding the context window of pre-trained LLMs from the conventional 4k up to over 2048k. As shown previously, our non-uniform positional interpolation method can yield an 8$\times$ increase in context length without the need for fine-tuning. However, when substantially larger extensions are desired (e.g., 512$\times$), fine-tuning becomes indispensable. A common strategy involves identifying suitable RoPE rescaling factors for the target 2048k length, followed by model fine-tuning. Nonetheless, this procedure is impeded by the significant computational resources required. Additionally, our observations indicate that fine-tuning LLMs for such large extension ratios is difficult to accomplish effectively (refer to Appendix).

Fortunately, LongRoPE demonstrates efficacy on both the base model and the extended, fine-tuned LLM. Consequently, we propose a streamlined, incremental approach that attains the desired 2048k context length using merely 1k fine-tuning steps, all confined within a training length of 256k.

$\Diamond$ \textit{LongRoPE search for pre-trained LLM extension to 256k}. Using LLaMA2 for illustration, we perform search experiments targeting context window sizes of 128k and 256k. At this point, the model's extension factors are 32$\times$ and 64$\times$, accordingly.

$\Diamond$ \textit{Fine-tuning to 256k}. Subsequently, we adapt the pre-trained LLM to support a context window of 256k. Our approach begins by fine-tuning LLaMA2 over 400 steps employing RoPE rescaled factors set to 128k. After this stage, we update the RoPE rescaled factors to 256k on the obtained checkpoint, followed by an extra 600 fine-tuning steps. This strategy has demonstrated greater efficiency compared to directly conducting fine-tuning for 256k.

$\Diamond$ \textit{Scaling fine-tuned extended LLM to 2048k using LongRoPE search}. In the final step, we carry out an additional search process on the LLM that has been fine-tuned for sequences of 256k length. This procedure successfully expands the context window to an impressive 2048k, with no need for supplementary fine-tuning. As a result, the context length is increased by a factor of $512\times$.

\textbf{Restoration within the original context window}. Upon scaling up to a 2048k-length context window, we observe diminished model effectiveness in the initial context span. This phenomenon aligns with established limitations of positional interpolation~\cite{pi}, wherein position embeddings from the original context space are compressed into a significantly smaller interval in higher dimensions, impairing the language model's capabilities. Specifically, applying a 512$\times$ scaling factor leads to severe congestion of position encodings within the first 4k tokens.

To address this, an additional evolution search is conducted on the expanded LLM to fine-tune the RoPE rescale factors specifically for shorter context lengths such as 4k and 8k. Because these shorter lengths necessitate less positional interpolation, we constrain the upper bound for the searched $\lambda$ accordingly. At inference time, the LLM adaptively modifies the relevant RoPE rescale factors.

\section{LongRoPE}

\label{sec:method}
Building on these insights, we propose LongRoPE, which initially incorporates a computationally efficient search algorithm designed to capitalize on both types of non-uniformity. Subsequently, this approach is leveraged to expand the LLM context window to exceed 2 million tokens.

\subsection{Problem Formulation}

These dual forms of non-uniformity may result in an expansive solution space and complicate the optimization process. To tackle this challenge, we recast the optimization of multidimensional non-uniform position interpolation as a search problem.

For a LLM targeting a  context window size of $L'$ and lengthy input documents $\mathbf{X}$, where each $\mathbf{x}\in\mathbf{X}$ surpasses $L'$ in token length, we denote the original rotary angle of the $i^{th}$ dimension in RoPE embedding at token position $n$ as  $\frac{n}{\beta^i}$. The optimization problem is then formulated as follows:

\begin{equation}
\begin{gathered}
		\label{eq:problem}

\underset {\mathbf{x}\in\mathbf{X}; \, |\mathbf{x}|\ge L'}{ \operatorname {arg\,min} } \,  \mathcal{L}\, \left( \text{LLM} (\text{RoPE}, \mathbf{X})\right),	\text{where \;} 
\\
\scalebox{0.93}{
$\underset {\begin{subarray}{l} i=0,\cdot\cdot,\frac{d}{2}-1;\\ n\in[ 0,|\mathbf{x}|);\end{subarray}}{\text{RoPE}_(n)}= {\left[\cdots,\cos\left(\mathbb{I}(\hat{\lambda}_{i}, \hat{n})\cdot\frac{n}{\beta^i}\right),\ \sin\left(\mathbb{I}(\hat{\lambda}_{i}, \hat{n})\cdot\frac{n}{\beta^i}\right),\cdots\right] }$}\\
	\text{where \;} \mathbb{I}(\hat{\lambda}_{i},\hat{n})=\begin{cases}
	1&\mbox{\;  $n<\hat{n}$}\\
{\frac{1}{\lambda}_{i}}&\mbox{\; $n\geq\hat{n}$}
\end{cases}
	\end{gathered}
\end{equation}

In this context, we define a group of rescale factors, $\mathbb{I}(\hat{\lambda}_{i},\hat{n})$, to address the two types of non-uniformity. Here, $\hat{\lambda_i}$ indicates the non-uniformity across RoPE dimensions, while $\hat{n}$ captures the variability in token positions. The function $\mathbb{I}(\hat{\lambda}_{i},\hat{n})$ serves to adjust the rotation angle for the $i^{th}$ dimension in RoPE; $\hat\lambda_{i}$ specifies the magnitude of rescaling, and $\hat{n}$ sets a threshold for token positions. For the initial $\hat{n}$-1 positions, rescaling does not occur and the standard RoPE rotation angle $\frac{n}{\beta^i}$ remains in effect. However, for token positions satisfying $n\ge\hat{n}$, the adjustment via the rescale factor is utilized.

Assuming a desired context window length of $L'$, our goal is to determine the best set of rescale factors ($\mathbb{I}(\hat{\lambda}_{0},\hat{n})$, $\mathbb{I}(\hat{\lambda}_{1},\hat{n})$, ...$\mathbb{I}(\hat{\lambda}_{i},\hat{n})$...) for each RoPE dimension from the $1^{st}$ through the $d^{th}$. By applying these rescaled $\text{RoPE}$ parameters to the target $\text{LLM}$, we aim to minimize the next token prediction loss, $\mathcal{L}$—or equivalently, the perplexity—when processing input sequences $\mathbf{X}$ of length $L'$.

\subsection{Searching the Non-uniform Position Interpolation}

To address the challenge outlined in Eq.~\ref{eq:problem}, we propose a straightforward but powerful approach that identifies optimal RoPE rescaling parameters, enabling comprehensive utilization of the multidimensional irregularities present in position embeddings.

\textbf{Search space}. We construct an extensive search space to accommodate the two types of non-uniformity. The details of this search space are provided in Table~\ref{tbl:searchspace}. In particular, our approach involves searching for a distinct rescale factor corresponding to each dimension within the RoPE embedding. To facilitate a more straightforward construction of the search space, we choose to optimize over $\lambda_i$ and $\hat{n}$, rather than directly searching for $\mathbb{I}(\hat{\lambda}_{i},\hat{n})$, where  $\hat{\lambda}_{i}=1/\lambda_i$. According to Table~\ref{tbl:searchspace}, the values for $\lambda_i$ are explored beginning from a lower bound of 1.0 (representing pure extrapolation) up to an upper bound of $s\times1.25$ (signifying greater interpolation than PI), with increments of 0.01. Here, $s$ denotes the ratio by which the target context window is extended.

$\hat{n}$ determines how many starting token positions maintain the original RoPE embedding, foregoing any position interpolation. In practice, we experiment with $\hat{n}$ values chosen from \{0, 1, 2, 4, 8, 12, 16, 20, 24, 28, 32, 64, 128, 256\}. Setting $\hat{n}=0$ means that every token position is assigned the computed rescale factors.

\textbf{Evolution-based search}. The search space outlined in Table~\ref{tbl:searchspace} encompasses a vast array of positional interpolation strategies, making exhaustive navigation computationally infeasible. As an illustration, scaling to a $s=4\times$ extension results in $400^{128/2}\times14 = 4\times10^{167}$ possible configurations. This combinatorial explosion becomes even more pronounced as the extension ratio increases. To cope with this, we adopt evolution search~\cite{guo2020single} and incorporate two targeted optimization methods that significantly enhance search efficiency. The comprehensive search workflow is depicted in Algorithm~\ref{alg:search}.

\textit{Optimized initial population generation}. In place of generating $P$ rescale factors through purely random initialization, we explicitly include the three RoPE rescale factors associated with PI, NTK, and YaRN among the first members of the population. The other $P-3$ individuals are produced by introducing random mutations to these three reference rescale factors, each with a mutation probability of $p$.

\textit{Monotonically non-decreasing constraint}. Once the initial population is produced, we evaluate each individual by calculating the LLM's perplexity. This involves applying the assigned RoPE rescale factors to the target LLM and then measuring the perplexity on input $\mathbf{X}$. The k highest-performing individuals are chosen as parents for the evolutionary process. Nonetheless, due to the extensive search space, uninformed mutation and crossover operations often lead to low-quality candidates, resulting in a surplus of redundant perplexity evaluations. This inefficiency becomes even more pronounced when $L'$ is large, because each perplexity inference is computationally expensive.

To tackle this challenge, we enforce a constraint of non-decreasing monotonicity on the selected RoPE rescaling coefficients: $\lambda_i \le \lambda_{i+1}$. Only those RoPE configurations adhering to this requirement are used in evaluating perplexity for the LLM, which greatly curtails the scope of the search. In particular, we stipulate that $\lambda_i$ must be monotonically increasing with respect to the RoPE dimension index ($i=0,\dots,63$). This choice is motivated by insights from NTK theory~\cite{jacot2018neural,tancik2020fourier,ntk}, which implies that dimensions associated with higher frequencies (i.e., lower indices) necessitate less interpolation (reflected in a smaller $\lambda_i$), whereas dimensions with lower frequencies (higher indices) permit greater interpolation (leading to a larger $\lambda_i$).

\begin{algorithm}[t]
	\small
	\caption{The search algorithm}
	\textbf{Input:} target LLM, input samples $\mathbf{X}$,   population size $P$, mutation size $N_1$, crossover size $N_2$, max iterations $\mathcal{T}$, mutate probability $p$\\

	\begin{algorithmic}[1]
		\label{alg:search}
		\STATE Top-k$\leftarrow\phi$;	
		\STATE $\text{P}_0 \leftarrow$ \textit{Initialize\_population\_with\_optimization} ($P$, $\mathbf{X}$, $p$); 
		\FOR{$i = 1$ to $\mathcal{T}$}
		\STATE \textit{Compute\_perplexity} (LLM, $\text{P}_{i-1}$, $\mathbf{X}$);
		\STATE  Top-k $\leftarrow$ \textit{Update\_Topk} (Top-k, $\text{P}_{i-1}$);
		\STATE $\text{P}_{mutation} \leftarrow$ \textit{Mutation\_with\_mono\_constraint} (Top-k, $N_1$, $p$); 
		\STATE $\text{P}_{crossover} \leftarrow$ \textit{Crossover\_with\_mono\_constraint} (Top-k, $N_2$);
		\STATE $\text{P}_i = \text{P}_{mutation} \cup \text{P}_{crossover} \cup $ Top-k;
		\ENDFOR
		\STATE Output the individual from Top-k exhibiting the minimal perplexity value;
	\end{algorithmic}
\end{algorithm}

\textbf{8$\times$ extension without fine-tuning}. Through evolutionary search, our approach successfully discovers non-uniform RoPE rescale factors that safeguard essential dimensions and positions, thereby reducing the loss of information that occurs during interpolation. As shown in Fig.\ref{fig:non-ft}, this technique enables us to expand the LLaMA2 context window from 4k up to 32k tokens without the need for fine-tuning. Conversely, previous approaches—including PI and non-uniform NTK and YaRN—exhibit sharp increases in perplexity when extending the context window beyond 2$\times$.

\subsection{Extending LLM Context Window to 2048K}
\label{sec:extendto2048k}

\textbf{Progressive extension to 2048k}. In this section, we present our strategy for expanding the context window of pre-trained LLMs from the standard 4k limit to more than 2048k tokens. As shown earlier, our non-uniform positional interpolation method is capable of delivering up to an 8$\times$ extension without any need for fine-tuning. However, when pushing toward significantly larger context windows (for instance, 512$\times$ extension), fine-tuning becomes indispensable. One approach involves identifying suitable RoPE rescaling factors tailored for the desired 2048k context length, followed by subsequent fine-tuning. This, however, comes with substantial obstacles given the extremely high demands on computational resources required for training. Furthermore, our practical observations indicate that effectively fine-tuning LLMs for such a significant increase in context window size is quite difficult (refer to Appendix).

Luckily, LongRoPE demonstrates efficacy when applied to both base and fine-tuned expanded LLMs. Consequently, we propose a streamlined, incremental approach that reaches the intended 2048k sequence length through only 1k fine-tuning steps, using training sequences of up to 256k in length.

$\Diamond$ \textit{Investigating LongRoPE search for scaling pre-trained LLMs to 256k}. Using LLaMA2 as a case study, we perform searches targeting expanded context window sizes of 128k and 256k tokens. This phase involves increasing the extension ratio to 32$\times$ and 64$\times$, correspondingly.

$\Diamond$ \textit{Fine-tuning to 256k}. Subsequently, we adapt the pre-trained LLM to extend its context window to 256k tokens. To do this, LLaMA2 is initially fine-tuned for 400 iterations utilizing RoPE rescaled factors set for 128k. After completing these steps, the RoPE rescaled factors are switched to 256k, and the model undergoes a further 600 steps of fine-tuning. This staged approach demonstrates greater efficiency compared to a single-phase fine-tuning directly targeting 256k.

$\Diamond$ \textit{Expanding fine-tuned extended LLM to 2048k using LongRoPE search}. As a final step, a secondary search is conducted utilizing the fine-tuned LLM trained at a 256k context length. This allows us to achieve a context window of 2048k—an exceptionally large size—without necessitating additional fine-tuning. The resulting total extension is a $512\times$ increase.

\textbf{Performance Degradation in Shorter Context Windows}.  Upon expanding the context window to a substantial length of 2048k, a reduction in performance is observed within the model's initial context window. This phenomenon is well-documented in studies of positional interpolation~\cite{pi}, where embedding positions at higher dimensions in the initial context range become confined to a smaller representational space, thereby impairing model effectiveness. Specifically, using a 512$\times$ expansion factor causes the positions inside the original 4k context window to cluster tightly, exacerbating the problem.

To address this, we conduct an additional evolution search on the extended LLM aimed at optimizing the RoPE rescale factors specifically for shorter context lengths, such as 4k and 8k. For these shorter sequences, we limit the maximum permissible $\lambda$ in the search space, reflecting the decreased need for positional interpolation. At inference time, the LLM automatically adapts the associated RoPE rescale factors in response to the context length.

 \section{Experiments}

\label{sec:exp}
\subsection{Setup}
\textbf{Evaluation Tasks and models}. LongRoPE is integrated into both LLaMA2-7B and Mistral-7B, with their effectiveness assessed across three dimensions: (1) the perplexity scores attained by long-context LLMs on extended-length documents; (2) the Passkey retrieval task, which tests the model’s capability to extract a specific key phrase amidst predominantly irrelevant content; and (3) standard benchmarks for LLMs, but constrained within a 4096-token context window.

\textbf{Fine-tuning}. For LLaMA2, we employ a learning rate of 2e-5 with a linear decay schedule and set the global batch size to 32. The model undergoes 400 fine-tuning steps on the Redpajama~\cite{redpajama} dataset, which is partitioned into 128k-length segments, each bounded by the BOS and EOS tokens. Subsequently, starting from the resulting checkpoint, we continue training for an additional 600 steps in order to reach a 256k context window. The 128k context variant is fine-tuned on 8 A100 GPUs utilizing the distributed training framework~\cite{cube}, whereas scaling to a 256k context window necessitates 16 A100 GPUs. For the Mistral model, we maintain a fixed learning rate of 1e-6, using a global batch size of 64. Both the 128k and 256k context models adhere to the configurations described in YaRN~\cite{yarn}: we train for 400 steps on Together Computer's Long-Data Collections~\cite{mistraldata}, each sequence having a length of 16k, and utilize 4 A100 GPUs throughout training.

\textbf{Search}. When the target window size does not exceed 256k, the following parameters are set: $P$=64, $N_1$=$N_2$=16, $p$=0.3, $\mathcal{T}$=40, with the top 32 candidates from each round chosen for mutation and crossover. Perplexity is assessed on 5 randomly selected PG19 validation samples, each at least as long as the required target context length. For window sizes larger than 512k, the numbers for population, mutation, and crossover are reduced by half. Perplexity evaluation in this scenario uses 3 randomly drawn samples from the Pile-Books3~\cite{pile} validation dataset.

\textbf{Baselines}. To extend the context length to 2048k, we conducted fine-tuning of models with 128k and 256k token context windows, resulting in LongRoPE-2048k (ft=128k) for LLaMA2 and LongRoPE-2048k (ft=256k) for Mistral. These four models are evaluated against leading baselines in context window extension, focusing on open-source LLMs that underwent fine-tuning following positional interpolation techniques such as PI, NTK, and YaRN. The comparison encompasses models like Together-32k~\cite{together}, Code LLaMA~\cite{codellama}, LongLoRA-full-FT-100k~\cite{longlora}, as well as YaRN-LLaMA and YaRN-Mistral~\cite{yarn}.

\subsection{Main Results}

\label{sec:mainresult}
\textbf{Language modeling for long sequences up to 256k tokens}. Our initial analysis focuses on benchmarking against leading long-context LLMs at an evaluation window of 256k. To illustrate the robustness of our approach, we employ two distinct datasets: the test sets from Proof-pile~\cite{pg19} and PG19~\cite{pile}. Perplexity measurements are obtained at multiple context sizes via a sliding window of 256 tokens. For PG19, our experiments encompass the entire test set comprising 100 documents. In the case of Proof-pile, adhering to the methodology of YaRN~\cite{yarn}, we randomly sample 10 instances, each exceeding a minimum length of 128k tokens.

Table~\ref{tbl:proof-pile} and Table~\ref{tbl:pg19} present a perplexity comparison between LLaMA2 and Mistral models expanded using various interpolation strategies on the Proof-pile and PG19 datasets, respectively. Two main findings emerge: \textbf{(1)} the perplexity of our extended models consistently decreases as the evaluation length increases from 4k up to 256k, indicating that these models effectively utilize longer contexts; \textbf{(2)} despite being extended to a context window 16$\times$ greater—a setting that usually hampers performance at shorter lengths—our LongRoPE-2048k models still achieve superior results over state-of-the-art baselines when evaluated within a 256k context length.

\textbf{Long sequence language modeling beyond 2000k}. To assess performance on very large-scale documents, we utilize the Books3~\cite{pile} dataset. For efficiency in evaluation, 20 books are randomly chosen, each with a length greater than 2048k, and a sliding window of 256k is applied.

As detailed in Table~\ref{tbl:books3}, LongRoPE is able to increase the context window for both LLaMA2-7B and Mistral-7B up to 2048k, while maintaining perplexity that is on par with or better than the baselines across shorter context windows ranging from 8k to 128k. Distinct performance trends emerge when comparing LLaMA2 and Mistral at the 2048k setting. Specifically, Mistral demonstrates superior performance relative to baselines for shorter sequences; however, its perplexity surpasses 7 when the context length exceeds 256k. For LLaMA2, results conform to expectations: perplexity continues to improve (i.e., decrease) with the increase in context length, apart from moderate upticks observed at 1024k and 2048k. Notably, for LLaMA2, LongRoPE-2048k achieves better results when the fine-tuning is performed at 256k rather than 128k, attributable to a lower secondary extension factor (8$\times$ versus 16$\times$). Conversely, Mistral yields improved results when fine-tuned using a window size of 128k. This discrepancy primarily stems from the adoption of a 16k training length, following the YaRN methodology, during fine-tuning at both 128k and 256k for Mistral. This choice impacts Mistral’s capability to further enlarge its context window post-fine-tuning.

\textbf{Passkey retrieval}. Next, we investigate how the size of the context window influences model performance in generation-based tasks. To do this, we utilize the synthetic passkey retrieval benchmark described in ~\cite{passkey}. In this setup, the model must locate and recover a randomly assigned five-digit passkey embedded within a lengthy document. The prompt template used for this evaluation can be found in the appendix. We conduct the passkey retrieval task 10 times, each with the passkey randomly inserted at a uniformly sampled position within the full evaluation context length.

Fig.~\ref{fig:passkey} presents a comparison of retrieval accuracy against baseline methods. The accuracy of current models declines steeply to zero once sequence length surpasses 128k tokens. In stark contrast, LongRoPE-LLaMA2-2048k (ft=256k) demonstrates robust performance in this difficult scenario, consistently achieving retrieval accuracy of at least 90\% across the entire range from 4k to 2048k tokens. Similarly, LongRoPE-Mistral-2048k (ft=128k) sustains perfect (100\%) retrieval up to 1800k tokens, with a decrease to 60\% accuracy at 2048k—a trend consistent with observations in Table~\ref{tbl:books3}, which shows a modest perplexity increase at the 2048k sequence length.

\textbf{Evaluation on standard benchmarks within the original context window}. We assess the performance of LongRoPE-2048k models using the Hugging Face Open LLM Leaderboard~\cite{huggingfaceboard} under both zero-shot and few-shot conditions. Specifically, the evaluation settings include 25-shot ARC-Challenge~\cite{allenai:arc}, 10-shot HellaSwag~\cite{zellers2019hellaswag}, 5-shot MMLU~\cite{mmlu}, and 0-shot TruthfulQA~\cite{lin2021truthfulqa}.

As indicated in Table~\ref{tbl:commonsense}, our models attain similar performance on the original benchmark—which was developed for shorter context windows—and even surpass the baseline Mistral on TruthfulQA by an additional 0.5\%. While LongRoPE-LLaMA2-2048k, after fine-tuning at 256k, exhibits a marginally greater decline in performance, its results largely stay within acceptable limits across the evaluated tasks.

\subsection{Ablation Results}

\textbf{Effectiveness of the second positional interpolation}. Within our progressive extension framework, we apply our search algorithm to perform a secondary, non-uniform positional interpolation on extended LLMs that have been fine-tuned. To assess the utility of this approach, we execute evaluations on our fine-tuned LLaMA2-256k model. This model is further scaled to 512k, 1024k, and 2048k positions using both PI and YaRN. As indicated in Table~\ref{tbl:secondsearch}, our method of non-uniform positional interpolation maintains perplexity at a stable level. Conversely, both PI and YaRN exhibit rapidly increasing perplexity as the extension ratio becomes larger.

\textbf{Effectiveness of recovery at shorter context lengths.} To address the degradation in model performance at shorter context windows, we fine-tune the RoPE scaling parameters for LongRoPE-2048k using our search strategy. In particular, we reduce the upper bound on the scale factors during the search process, aiming to limit excessive interpolation for 4k and 8k context lengths. Table~\ref{tbl:shortperformance} presents a perplexity analysis of LongRoPE-LLaMA2-2048k on the Proof-pile dataset for both 4k and 8k contexts, along with corresponding average LLM benchmark scores. These findings illustrate a notable boost in performance at reduced context lengths.

\textbf{Analysis on the two forms of non-uniformities}. In this section, we conduct ablation studies to disentangle the impact of each non-uniformity component on model performance. Two separate experiments are designed: (i) expanding LLaMA2-7B to accommodate sequence lengths of 16k and 32k via multiple approaches—PI, optimizing only the RoPE dimension, and optimizing both identified non-uniformities; (ii) scaling our fine-tuned LLaMA2 (256k-length) model up to 2048k using analogous strategies. Perplexity is measured prior to any additional fine-tuning. According to Table~\ref{tbl:searchdimension}, introducing non-uniformity in the RoPE dimension yields a notable decrease in perplexity over linear interpolation as used in PI. Adjusting for token position non-uniformity further enhances results for 16k and 32k input lengths, although its benefit diminishes at 2048k, likely attributable to the extreme length regime. Merely retaining initial tokens, forgoing interpolation, appears ineffective for such long sequences; exploring this topic will be left to subsequent investigations.

\section{Related Works}

Besides position interpolation techniques, this section reviews other related methodologies.

\textbf{Retrieval-based} approaches utilize external memory components to store information from earlier contexts, and employ retrieval mechanisms to fetch relevant documents during inference~\cite{longllama,wang2023augmenting,borgeaud2022improving}. Generally, these systems require substantial architectural changes to the LLMs. By comparison, our method introduces only minimal changes to positional embeddings, making it significantly more lightweight. Furthermore, our approach extends to a wider range of long context applications beyond retrieval, such as long document summarization and few-shot learning.

\textbf{Attention-based context window extensions}. In addition to interpolating positional embeddings, other studies have managed to extend the effective input context by adjusting attention mechanisms while keeping the original LLM context window unchanged~\cite{han2023lminfinite, streamingllm,parallelwindow}. The principal strategy involves employing innovative attention masks to address the attention explosion problem introduced by incorporating new positions. These approaches are complementary to positional interpolation techniques.

\textbf{Fine-tuning based approaches} investigate strategies for adapting pre-trained LLMs to longer input sequences by altering position embeddings. Examples such as Code LLaMA~\cite{codellama}, LLaMA2 Long~\cite{xiong2023effective}, and ScaledRoPE~\cite{liu2023scaling} adopt significantly larger base values for RoPE before fine-tuning on the desired input lengths. In contrast, our approach is adaptable to a variety of target lengths and is capable of handling sequences extending beyond 2M tokens. Additionally, given that fine-tuning for extended context lengths (i.e., surpassing 128k) can be prohibitively resource-intensive, solutions like LongLoRA~\cite{longlora} and PoSE~\cite{zhu2023pose} have been introduced to reduce the associated computational burden. Our method remains complementary to these efficiency-focused fine-tuning approaches.

\section{Conclusion}

This paper introduces LongRoPE, a technique that significantly enhances the context length of LLMs up to 2048k, while preserving their effectiveness on tasks within the original, shorter context window. Our approach utilizes two distinct types of non-uniformity in RoPE positional encoding, identified through a streamlined evolutionary search process. This methodology offers two main advantages: it yields strong initial weights for subsequent fine-tuning and supports an 8$\times$ increase in the context window size even without fine-tuning. Furthermore, we develop a progressive expansion scheme that starts with LLMs fine-tuned on a 256k context and incrementally reaches a 2048k context window, all without necessitating additional rounds of fine-tuning. Our comprehensive empirical results demonstrate the efficacy of LongRoPE. We anticipate that the LongRoPE-2048k models will open up a wide range of novel long-context applications and stimulate future investigations in this area.

\section*{Broader Impacts} This study aims to contribute to progress within the Machine Learning community. While various societal impacts could arise from our research, we do not identify any particular effects that warrant explicit mention at this time.

	\balance

\end{document}